\documentclass{article}
\usepackage{amsmath,amssymb,amsthm}
\usepackage{algorithm}
\usepackage{algpseudocode}
\usepackage{enumitem}
\usepackage{hyperref}
\newtheorem{lemma}{Lemma}
\newtheorem{theorem}{Theorem}
\newtheorem{corollary}{Corollary}
\allowdisplaybreaks

\begin{document}

\section*{Proximal Gradient Method (Forward–Backward Splitting)}

\section*{Mathematical Formulation}
Let 
\[
F(x)\;:=\;f(x)+g(x),\qquad x\in\mathbb{R}^{d},
\]
where  

\begin{itemize}[nosep]
\item $f:\mathbb{R}^{d}\to\mathbb{R}$ is convex and differentiable with $L$–Lipschitz continuous gradient,
      \[
      \|\nabla f(u)-\nabla f(v)\|\le L\|u-v\|\quad\forall u,v\in\mathbb{R}^{d};
      \]
\item $g:\mathbb{R}^{d}\to\mathbb{R}\cup\{+\infty\}$ is proper, closed and convex (possibly non‑smooth).
\end{itemize}
For a stepsize $\eta>0$ the \emph{proximal operator} of $g$ is
\[
\operatorname{prox}_{\eta g}(z)\;:=\;\arg\min_{x\in\mathbb{R}^{d}}\Bigl\{ g(x)+\frac{1}{2\eta}\|x-z\|^{2}\Bigr\}.
\]

The proximal gradient iteration is
\begin{equation}\label{eq:pg}
x_{k+1}\;=\;\operatorname{prox}_{\eta g}\!\bigl(x_{k}-\eta\nabla f(x_{k})\bigr),\qquad k=0,1,2,\dots
\end{equation}
with a fixed stepsize $\eta\in(0,1/L]$ unless stated otherwise.

\section*{Pseudocode}
\begin{algorithm}[H]
\caption{Proximal Gradient Method}
\begin{algorithmic}[1]
\Require Initial point $x_{0}\in\mathbb{R}^{d}$, stepsize $\eta\in(0,1/L]$, maximum iterations $K$.
\For{$k=0$ \textbf{to} $K-1$}
    \State $g_{k}\gets \nabla f(x_{k})$                         \Comment{Gradient of the smooth part}
    \State $y_{k}\gets x_{k}-\eta\,g_{k}$                       \Comment{Forward step}
    \State $x_{k+1}\gets \operatorname{prox}_{\eta g}(y_{k})$    \Comment{Backward (prox) step}
\EndFor
\State \Return $x_{K}$
\end{algorithmic}
\end{algorithm}

\section*{Dry Run Example}
Consider the scalar problem
\[
f(w)=(w-3)^{2},\qquad g\equiv0,
\]
so that $\operatorname{prox}_{\eta g}(z)=z$.  
Take $w_{0}=0$, stepsize $\eta=0.1$, and run three iterations.

\begin{center}
\begin{tabular}{c|c|c|c}
$k$ & $w_{k}$ & $\nabla f(w_{k})$ & $w_{k+1}=w_{k}-\eta\nabla f(w_{k})$\\ \hline
0 & $0$ & $2(w_{0}-3)=-6$ & $0-0.1(-6)=0.6$\\
1 & $0.6$ & $2(0.6-3)=-4.8$ & $0.6-0.1(-4.8)=1.08$\\
2 & $1.08$ & $2(1.08-3)=-3.84$ & $1.08-0.1(-3.84)=1.464$\\
\end{tabular}
\end{center}
The objective values are $F(w_{0})=9$, $F(w_{1})=(0.6-3)^{2}=5.76$, $F(w_{2})= (1.08-3)^{2}=3.6864$, $F(w_{3})=(1.464-3)^{2}=2.3630$, showing monotone decrease.

\newpage
\section*{Assumptions}
\begin{enumerate}[label={\bf(A\arabic*)}, leftmargin=*]
\item \label{A1} \textbf{Smoothness.} $f$ is convex and differentiable with $L$‑Lipschitz gradient:
\[
\|\nabla f(u)-\nabla f(v)\|\le L\|u-v\|,\qquad\forall u,v\in\mathbb{R}^{d}.
\]

\item \label{A2} \textbf{Convexity of $g$.} $g$ is proper, closed and convex (no smoothness required).

\item \label{A3} \textbf{Stepsize.} The constant stepsize satisfies
\[
0<\eta\le \frac{1}{L}.
\]

\item \label{A4} \textbf{Existence of a minimizer.} The set 
\[
\mathcal{X}^{\star}:=\arg\min_{x}F(x)
\]
is non‑empty and we denote any optimal point by $x^{\star}$.
\end{enumerate}

\section*{Main Convergence Theorem}
\begin{theorem}\label{thm:rate}
Under Assumptions \ref{A1}–\ref{A4}, the iterates generated by \eqref{eq:pg} satisfy for every $k\ge 0$
\begin{equation}\label{eq:func-decrease}
F(x_{k})-F(x^{\star})\;\le\;\frac{\|x_{0}-x^{\star}\|^{2}}{2\eta\,k}.
\end{equation}
Consequently $F(x_{k})\to F(x^{\star})$ and the convergence rate in function value is $O(1/k)$.
\end{theorem}

\section*{Theoretical Properties}
\begin{itemize}[nosep]
\item \textbf{Stability.} The iteration is a non‑expansive map composed of a gradient step (Lipschitz) and a firmly non‑expansive proximal map; hence the sequence $\{x_{k}\}$ remains bounded whenever a minimizer exists.

\item \textbf{Hyperparameter Sensitivity.} The stepsize bound $\eta\le 1/L$ is tight: for $\eta>1/L$ the descent lemma (Lemma~\ref{lem:descent}) may fail, and the method can diverge.

\item \textbf{Comparison to Gradient Descent.} When $g\equiv0$, $\operatorname{prox}_{\eta g}$ reduces to the identity and \eqref{eq:pg} becomes classic gradient descent with the same $O(1/k)$ rate.

\item \textbf{Special case $g$ simple.} If $g$ admits a closed form proximal operator (e.g., $\ell_{1}$ norm), each iteration is inexpensive, making the method suitable for large‑scale composite optimization.
\end{itemize}

\section*{Discussion}
The $O(1/k)$ bound matches the optimal rate for first‑order methods on the class of convex, Lipschitz‑smooth problems. Acceleration (Nesterov’s momentum) can improve the rate to $O(1/k^{2})$, but the basic proximal gradient method already enjoys a simple implementation and strong theoretical guarantees derived solely from convexity and the non‑expansiveness of the proximal map.

\newpage
\section*{Preliminaries (Definitions and Lemmas)}

\begin{lemma}[Firm non‑expansiveness of the proximal operator]\label{lem:firm}
For any $\eta>0$ and any $u,v\in\mathbb{R}^{d}$,
\[
\big\|\operatorname{prox}_{\eta g}(u)-\operatorname{prox}_{\eta g}(v)\big\|^{2}
\le \langle u-v,\;\operatorname{prox}_{\eta g}(u)-\operatorname{prox}_{\eta g}(v)\rangle .
\]
\end{lemma}
\begin{proof}
Standard result; follows from the optimality condition of the proximal sub‑problem and monotonicity of $\partial g$. \qedhere
\end{proof}

\begin{lemma}[Descent lemma for $L$‑smooth $f$]\label{lem:descent}
If $\nabla f$ is $L$‑Lipschitz, then for any $u,v\in\mathbb{R}^{d}$,
\[
f(v)\le f(u)+\langle\nabla f(u),v-u\rangle+\frac{L}{2}\|v-u\|^{2}.
\]
\end{lemma}
\begin{proof}
Integrate the Lipschitz condition on the gradient; see e.g. \cite{Nesterov2004}. \qedhere
\end{proof}

\section*{Main Proof (Step‑by‑Step Derivation)}
We prove Theorem~\ref{thm:rate}. Throughout the proof we fix an arbitrary optimal point $x^{\star}\in\mathcal{X}^{\star}$.

\noindent\textbf{Notation.} Set 
\[
y_{k}=x_{k}-\eta\nabla f(x_{k}),\qquad x_{k+1}= \operatorname{prox}_{\eta g}(y_{k}).
\]

\begin{enumerate}
\item[Step 1.] \textbf{Optimality condition of the proximal step.}  
By definition of $\operatorname{prox}_{\eta g}$,
\[
0\in \partial g(x_{k+1})+\frac{1}{\eta}(x_{k+1}-y_{k})
\quad\Longrightarrow\quad
-\frac{1}{\eta}(x_{k+1}-y_{k})\in\partial g(x_{k+1}). \tag{1}
\]

\item[Step 2.] \textbf{Convexity of $g$.}  
For any $z$ and any $s\in\partial g(x_{k+1})$,
\[
g(z)\ge g(x_{k+1})+\langle s,\,z-x_{k+1}\rangle .
\]
Choosing $z=x^{\star}$ and using (1) yields
\[
g(x^{\star})\ge g(x_{k+1})-\frac{1}{\eta}\langle x_{k+1}-y_{k},\,x^{\star}-x_{k+1}\rangle . \tag{2}
\]

\item[Step 3.] \textbf{Apply Lemma~\ref{lem:descent} to $f$.}  
With $u=x_{k}$ and $v=x_{k+1}$,
\[
f(x_{k+1})\le f(x_{k})+\langle\nabla f(x_{k}),x_{k+1}-x_{k}\rangle+\frac{L}{2}\|x_{k+1}-x_{k}\|^{2}. \tag{3}
\]

\item[Step 4.] \textbf{Combine $f$ and $g$ inequalities.}  
Add (2) and (3):
\[
\begin{aligned}
F(x_{k+1}) &\le f(x_{k})+g(x_{k+1})
            +\langle\nabla f(x_{k}),x_{k+1}-x_{k}\rangle
            +\frac{L}{2}\|x_{k+1}-x_{k}\|^{2} \\
           &\quad -\frac{1}{\eta}\langle x_{k+1}-y_{k},\,x^{\star}-x_{k+1}\rangle .
\end{aligned}
\tag{4}
\]

\item[Step 5.] \textbf{Express $y_{k}$ in terms of $x_{k}$.}  
Recall $y_{k}=x_{k}-\eta\nabla f(x_{k})$, hence
\[
x_{k+1}-y_{k}=x_{k+1}-x_{k}+\eta\nabla f(x_{k}). \tag{5}
\]

\item[Step 6.] \textbf{Insert (5) into the inner product of (4).}  
\[
\begin{aligned}
-\frac{1}{\eta}\langle x_{k+1}-y_{k},\,x^{\star}-x_{k+1}\rangle
&= -\frac{1}{\eta}\langle x_{k+1}-x_{k},\,x^{\star}-x_{k+1}\rangle
   -\langle\nabla f(x_{k}),\,x^{\star}-x_{k+1}\rangle .
\end{aligned}
\tag{6}
\]

\item[Step 7.] \textbf{Plug (6) into (4) and cancel the gradient terms.}  
The terms $+\langle\nabla f(x_{k}),x_{k+1}-x_{k}\rangle$ (from (4)) and $-\langle\nabla f(x_{k}),x^{\star}-x_{k+1}\rangle$ (from (6)) combine to give
\[
\langle\nabla f(x_{k}),\,x_{k+1}-x_{k}+x_{k+1}-x^{\star}\rangle
= \langle\nabla f(x_{k}),\,x_{k+1}-x^{\star}\rangle .
\]
Thus (4) becomes
\[
F(x_{k+1})\le f(x_{k})+g(x_{k+1})
            -\frac{1}{\eta}\langle x_{k+1}-x_{k},\,x^{\star}-x_{k+1}\rangle
            +\frac{L}{2}\|x_{k+1}-x_{k}\|^{2}. \tag{7}
\]

\item[Step 8.] \textbf{Upper bound the mixed inner product.}  
Using the identity $2\langle a,b\rangle = \|a\|^{2}+\|b\|^{2}-\|a-b\|^{2}$ with
$a=x_{k+1}-x_{k}$ and $b=x^{\star}-x_{k+1}$,
\[
-\frac{1}{\eta}\langle x_{k+1}-x_{k},\,x^{\star}-x_{k+1}\rangle
= \frac{1}{2\eta}\Bigl(\|x_{k}-x^{\star}\|^{2}
                     -\|x_{k+1}-x^{\star}\|^{2}
                     -\|x_{k+1}-x_{k}\|^{2}\Bigr). \tag{8}
\]

\item[Step 9.] \textbf{Insert (8) into (7).}  
\[
\begin{aligned}
F(x_{k+1}) &\le f(x_{k})+g(x_{k+1})
            +\frac{1}{2\eta}\bigl(\|x_{k}-x^{\star}\|^{2}
                                 -\|x_{k+1}-x^{\star}\|^{2}
                                 -\|x_{k+1}-x_{k}\|^{2}\bigr)
            +\frac{L}{2}\|x_{k+1}-x_{k}\|^{2}.
\end{aligned}
\tag{9}
\]

\item[Step 10.] \textbf{Replace $f(x_{k})$ by $F(x_{k})-g(x_{k})$ and use convexity of $g$.}  
Since $g$ is convex,
\[
g(x_{k})\ge g(x_{k+1})+\langle s,\,x_{k}-x_{k+1}\rangle
\]
for some $s\in\partial g(x_{k+1})$. From the optimality condition (1) we have $s=-\frac{1}{\eta}(x_{k+1}-y_{k})$, which yields $g(x_{k})\ge g(x_{k+1})-\frac{1}{\eta}\langle x_{k+1}-y_{k},\,x_{k}-x_{k+1}\rangle$. This term is non‑negative because of the firm non‑expansiveness (Lemma~\ref{lem:firm}). Consequently,
\[
g(x_{k})\ge g(x_{k+1}) .
\]
Hence $f(x_{k})+g(x_{k+1})\le F(x_{k})$ and (9) simplifies to
\[
F(x_{k+1})\le F(x_{k})
            +\frac{1}{2\eta}\bigl(\|x_{k}-x^{\star}\|^{2}
                                 -\|x_{k+1}-x^{\star}\|^{2}
                                 -\|x_{k+1}-x_{k}\|^{2}\bigr)
            +\frac{L}{2}\|x_{k+1}-x_{k}\|^{2}. \tag{10}
\]

\item[Step 11.] \textbf{Use the stepsize bound $\eta\le 1/L$.}  
Since $\eta\le 1/L$, we have $L\le 1/\eta$, thus
\[
\frac{L}{2}\|x_{k+1}-x_{k}\|^{2}
\le \frac{1}{2\eta}\|x_{k+1}-x_{k}\|^{2}.
\]
Insert into (10) to cancel the $\|x_{k+1}-x_{k}\|^{2}$ terms:
\[
F(x_{k+1})\le F(x_{k})
            +\frac{1}{2\eta}\bigl(\|x_{k}-x^{\star}\|^{2}
                                 -\|x_{k+1}-x^{\star}\|^{2}\bigr). \tag{11}
\]

\item[Step 12.] \textbf{Rearrange (11).}  
\[
\frac{1}{2\eta}\bigl(\|x_{k}-x^{\star}\|^{2}
                     -\|x_{k+1}-x^{\star}\|^{2}\bigr)
\ge F(x_{k+1})-F(x_{k}). \tag{12}
\]

\item[Step 13.] \textbf{Telescoping sum.}  
Sum (12) for $k=0,\dots, K-1$:
\[
\frac{1}{2\eta}\bigl(\|x_{0}-x^{\star}\|^{2}
                     -\|x_{K}-x^{\star}\|^{2}\bigr)
\ge \sum_{k=0}^{K-1}\bigl(F(x_{k+1})-F(x_{k})\bigr)
=F(x_{K})-F(x_{0}). \tag{13}
\]

\item[Step 14.] \textbf{Upper bound the left‑hand side by discarding the non‑negative term $\|x_{K}-x^{\star}\|^{2}$.}
\[
\frac{1}{2\eta}\|x_{0}-x^{\star}\|^{2}\ge F(x_{K})-F(x_{0}). \tag{14}
\]

\item[Step 15.] \textbf{Express the bound for any $k\ge 1$.}  
From (11) we also have
\[
F(x_{k})-F(x^{\star})\le\frac{1}{2\eta}\bigl(\|x_{k-1}-x^{\star}\|^{2}
                                 -\|x_{k}-x^{\star}\|^{2}\bigr).
\]
Summing the inequality for $k=1,\dots,K$ and dividing by $K$ yields
\[
F(x_{K})-F(x^{\star})\le\frac{\|x_{0}-x^{\star}\|^{2}}{2\eta K}. \tag{15}
\]

\end{enumerate}
Equation~(15) is exactly the claim \eqref{eq:func-decrease} of Theorem~\ref{thm:rate}. ∎

\section*{Rate Derivation (With Constants)}
From \eqref{eq:func-decrease} we obtain the explicit constant
\[
C\;:=\;\frac{\|x_{0}-x^{\star}\|^{2}}{2\eta},
\qquad\text{so that}\qquad
F(x_{k})-F(x^{\star})\le \frac{C}{k},\;k\ge1.
\]
If $g\equiv0$, $C$ reduces to the usual gradient‑descent constant. For $\eta=1/L$ the bound becomes
\[
F(x_{k})-F(x^{\star})\le \frac{L\|x_{0}-x^{\star}\|^{2}}{2k}.
\]

\section*{Special Cases}
\begin{enumerate}[label={\bf(S\arabic*)}]
\item \textbf{Pure smooth case $g\equiv0$.} The proximal step is identity and the method coincides with gradient descent; the proof simplifies by dropping Lemma~\ref{lem:firm}.

\item \textbf{Indicator function $g=\iota_{\mathcal{C}}$ of a closed convex set $\mathcal{C}$.} Then $\operatorname{prox}_{\eta g}$ is the Euclidean projection onto $\mathcal{C}$. The algorithm becomes projected gradient descent with the same $O(1/k)$ guarantee.

\item \textbf{$\ell_{1}$ regularisation $g(x)=\lambda\|x\|_{1}$.} The proximal operator is soft‑thresholding; the theorem guarantees convergence of the composite objective (LASSO) at rate $O(1/k)$ without any extra assumptions beyond convexity.

\end{enumerate}

\begin{thebibliography}{9}
\bibitem{Nesterov2004}
Y.~Nesterov, \emph{Introductory Lectures on Convex Optimization: A Basic Course}, Kluwer Academic Publishers, 2004.
\end{thebibliography}

\end{document}